\chapter{What pii is}

\label{ch:what}

pii is a \emph{personal information index}: one Rust binary that keeps stores of runs,
memory, entities and lessons on disk, executes \emph{graphs} --- small programs written as
data --- on a daemon, and records every state change in a journal that can be replayed,
counted and audited. This chapter answers three questions before the rest of the manual
asks anything of you: what the program does, what it is made of, and which of its answers
you are allowed to trust. Every number and path in this chapter was read from the
repository at the revision the manual was built from, not from a design document; where a
mechanism exists because something failed, the failure is named.

It is not a chat application that happens to have a database. It is an execution engine
whose every action leaves evidence, and the questions you will ask of it are of the form
\enquote{what ran, what did it decide, and can I prove it}. The design follows from that:
if a fact cannot be shown from the journal or from a store, pii treats it as unknown rather
than as true. The three exit codes of \texttt{pii health} are the smallest example --- a
daemon that answers from a different build than your client is reported as drift, not as
health.

\section{The three things it does}

\begin{enumerate}
	\item \textbf{Runs work.} A graph is a set of nodes --- ask a model, run a command,
	judge a set of exit criteria, call another graph --- connected by named ports. pii
	executes it, journals each transition, and returns a verdict rather than a transcript.
	\item \textbf{Remembers what it learned.} Runs, node-call events, memory documents,
	entities and lessons live in redb stores with counters, snapshots and, since the
	retention row, an age-based policy that bounds them.
	\item \textbf{Tells the truth about itself.} \texttt{pii health} answers from a record
	kept current by the write path rather than by replaying history at read time, and the
	answer names the build, the config file's hash, each store's path and format version,
	and the mounts the sandbox actually granted.
\end{enumerate}

Those three are one mechanism seen from three sides. The journal is what makes a verdict
provable; the stores are what make a second run cheaper than the first; the identity
payload is what makes an upgrade, a rollback or a wrong state directory visible instead of
merely suspected.

\section{The shape of the system}

The workspace is twenty-two \texttt{pii-*} crates plus \texttt{xtask}, the internal
tool that runs the quality gate. Four binaries come out of it: \texttt{pii} (the daemon
and the command line in one file), \texttt{pii-mcp} (a Model Context Protocol server),
\texttt{pii-web} (a small HTTP surface) and \texttt{pii-evolve}. Only \texttt{pii} is
required to operate the product; the installer places the other three beside it, and
\texttt{pii-mcp} is staged as \texttt{pii-mcp.new} rather than swapped in, because its
callers are other programs whose behaviour has to be measured before the swap.

\begin{longtable}{@{}p{0.20\textwidth}p{0.30\textwidth}p{0.42\textwidth}@{}}
	\toprule
	\textbf{Binary} & \textbf{Crate} & \textbf{What it is for} \\
	\midrule
	\endhead
	\texttt{pii} & \texttt{pii-cli} & The command line, and the daemon when invoked as
	\texttt{pii daemon}. Embeds the engine in-process for a one-shot prompt. \\
	\texttt{pii-mcp} & \texttt{pii-mcp} & A Model Context Protocol server over the same
	engine, for tools that speak MCP. Installed as \texttt{.new}; the swap is an operator
	step. \\
	\texttt{pii-web} & \texttt{pii-web} & A separate HTTP binary, not linked into
	\texttt{pii}. \\
	\texttt{pii-evolve} & \texttt{pii-evolve} & A separate evolution binary, not linked
	into \texttt{pii}. \\
	\bottomrule
\end{longtable}

Inside \texttt{pii} the layering is strict, and it is the reason an operator can reason
about failure at all. The control commands --- \texttt{run}, \texttt{queue},
\texttt{health}, \texttt{memory}, \texttt{skill} --- never open a store: they dial
\texttt{control.sock}. The daemon owns the stores, the run pool, the timer wheel and the
channels. The engine executes one run at a time per worker and journals every node
transition. \texttt{pii-store} is the only crate that knows redb exists, and the commands
that do open a file (\texttt{store}, \texttt{backup}, \texttt{restore}, and the offline
\texttt{graph run} and \texttt{graph trace} with \texttt{--db}) take redb's exclusive lock
and refuse while a daemon holds it.

\begin{center}
\resizebox{\linewidth}{!}{
\begin{tikzpicture}[node distance=9mm and 14mm]
  \node[b] (cli) {\texttt{pii} command line\\\scriptsize one subcommand per invocation};
  \node[b, right=of cli] (daemon) {daemon\\\scriptsize owns every store};
  \node[b, right=of daemon] (engine) {engine\\\scriptsize one run per worker};
  \node[s, below=of engine] (redb) {redb files\\\scriptsize seven stores, two in durable/};
  \node[b, below=of daemon] (pool) {run pool\\\scriptsize bounded threads};
  \node[b, below=of cli] (oneshot) {one-shot\\\scriptsize engine in-process};
  \draw[f] (cli) -- node[p,above=2pt,midway]{control.sock} (daemon);
  \draw[f] (daemon) -- (engine);
  \draw[f] (engine) -- (redb);
  \draw[sf] (daemon) -- (pool);
  \draw[sf] (cli) to[bend right=14] node[p,left=2pt,midway]{no daemon?} (oneshot);
  \draw[sf] (oneshot) -- (engine);
\end{tikzpicture}}
\captionof{figure}{Two ways in and one owner. The command line either reaches the daemon over
its socket, or --- for a one-shot prompt only --- builds the engine in its own process. Notice
that the store has exactly one writer, which is what makes an offline command refuse rather than
contend.}
\end{center}

The distinction matters in practice. \texttt{pii store compact} and \texttt{pii backup}
are two of the commands that do open the files, and both refuse while a daemon holds them
instead of copying a store mid-write. Everything else --- \texttt{run}, \texttt{queue},
\texttt{health}, \texttt{memory}, \texttt{skill} --- is a protocol call.

\section{A graph is data}

A graph file is TOML with five top-level facts: a \texttt{name} (the key \texttt{pii run
start} resolves), an \texttt{entry} node, an optional \texttt{exit} node, a \texttt{caps}
list, and an optional \texttt{[limits]} table. Nodes and edges follow. This is the
beginning of the development queue's own graph, at the revision this manual describes:

\begin{lstlisting}[language=bash]
name = "queue"
entry = "select"
exit = "record"
caps = ["spawn", "proc", "fs_read", "fs_write"]

[limits]
max_wall_ms = 6000000

[[node]]
id = "select"
kind = "select"

[[node]]
id = "implement"
kind = "call"
graph = "spec"
sha256 = "771d098b39db7369dd71ed1fb23ffa63b90b25fed7788c830e4bfb214c424443"
caps = ["spawn", "proc", "fs_read", "fs_write"]

[[edge]]
from = "select"
port = "ok"
to = "implement"
kind = "data"
\end{lstlisting}

Three things in those lines carry most of the design. The \texttt{sha256} on a
\texttt{call} is a pin: the callee is spliced into the host at load time, so a change to
the callee invalidates the pin until someone re-pins it deliberately. The \texttt{caps} on
the call node are the host's promise that it holds what the callee needs. And the
\texttt{port} on an edge is a name, not a boolean --- \texttt{select} leaves by
\texttt{ok}, \texttt{empty} or \texttt{blocked}, and each can go somewhere different.

\begin{center}
\resizebox{\linewidth}{!}{
\begin{tikzpicture}[node distance=8mm and 12mm]
  \node[b] (entry) {entry node\\\scriptsize \texttt{select}};
  \node[v, right=of entry] (port) {port?};
  \node[b, right=of port, yshift=10mm] (ok) {\texttt{implement}\\\scriptsize a call, pinned};
  \node[b, right=of port, yshift=-10mm] (empty) {\texttt{record}\\\scriptsize nothing to do};
  \node[b, below=of entry] (caps) {\texttt{caps}\\\scriptsize the ceiling for every node};
  \node[b, right=of ok] (exit) {exit node\\\scriptsize \texttt{record}};
  \draw[f] (entry) -- (port);
  \draw[f] (port) -- node[p,above=2pt,midway]{ok} (ok);
  \draw[f] (port) -- node[p,above=2pt,midway]{empty} (empty);
  \draw[f] (ok) -- (exit);
  \draw[bk] (caps) -- (entry);
\end{tikzpicture}}
\captionof{figure}{A graph from the operator's side: one entry, named ports, and an exit that
makes the graph usable as a part of another graph. A graph without an exit is a run, not a
component --- nothing can call it.}
\end{center}

The node set is small by design. Seven kinds live in \texttt{pii-graph} and need no
services at all: \texttt{transform}, \texttt{route}, \texttt{emit}, \texttt{wait},
\texttt{timer}, \texttt{refuse} and \texttt{require}. The daemon registers the rest,
because they need a store, a socket or a process.

\begin{longtable}{@{}p{0.17\textwidth}p{0.75\textwidth}@{}}
	\toprule
	\textbf{Kind} & \textbf{What it does, and what it costs} \\
	\midrule
	\endhead
	\texttt{transform} & Rewrites the value on the port. No process, no model. \\
	\texttt{route} & Chooses a port from the value it received. The cheapest decision in
	the system. \\
	\texttt{wait}, \texttt{timer} & Sleeps, or waits for a wall-clock moment. \texttt{timer}
	is what makes a graph schedulable. \\
	\texttt{refuse}, \texttt{require} & Ends the run with a named refusal, or refuses to
	start unless a condition holds. A refusal is a verdict, not an exception. \\
	\texttt{coder} & Spawns a coding agent (\texttt{pi} or \texttt{claude}) and takes the
	agent's verdict from the run of its criteria, never from the agent's own report. \\
	\texttt{judge} & Runs a spec's exit criteria through a shell and reports
	\texttt{ok}, \texttt{fail} or \texttt{unrunnable}. \\
	\texttt{select}, \texttt{record} & Claim one work item from the queue store; write a
	verdict back to it. \\
	\texttt{agent}, \texttt{tool}, \texttt{worker} & Call a model, a built-in tool, or a
	supervised child process. All three need \texttt{spawn} or \texttt{proc}. \\
	\texttt{gate} & Suspends the run for an approval, and resumes it on the decision. \\
	\texttt{post}, \texttt{mqtt} & Send a message to a configured channel, or publish to a
	broker. \\
	\texttt{memory}, \texttt{remember}, \texttt{recall} & Consolidate the memory store;
	write a fact from a run; read it back in a later run. \\
	\texttt{heal\_detect}, \texttt{heal\_record} & Run a gate, record the failures it
	found, and keep the record. \\
	\texttt{spawn} & Starts a child run on the same engine. \\
	\texttt{call} & Splices another graph in at load time. Not a runtime kind: it becomes
	the callee's nodes before the run starts. \\
	\bottomrule
\end{longtable}

A graph that lints can still be wrong. A graph that does not lint will not run:
\texttt{pii graph lint} is the only check that refuses an edge to a port the node does not
declare, a capability the graph does not hold, or a pin that no longer matches the
callee's file.

\section{The journal, and the verdicts}

Every run writes an append-only journal, and the events are a fixed vocabulary of
seventeen: \texttt{Started}, \texttt{TickStarted}, \texttt{NodeStarted},
\texttt{NodeFinished} (with its port and the hash of the value), \texttt{NodeFailed},
\texttt{NodeProgress}, \texttt{NodeCancelled}, \texttt{NodeDead}, \texttt{NodeInterrupted},
\texttt{Suspended}, \texttt{Resumed}, \texttt{NodeObserved}, \texttt{Finished},
\texttt{EffectStarted}, \texttt{EffectConfirmed}, \texttt{NodeRetried} and
\texttt{RunSpawned}. That vocabulary is the reason a trace can be read rather than guessed
at. Two of the seventeen exist only for effects: an effect node journals its idempotency
key \emph{before} it acts, so a key with no confirmation means the effect may or may not
have happened, and a replay says so instead of acting twice. A run's stored state is one
of five:

\begin{center}
\begin{tabular}{@{}p{0.16\textwidth}p{0.70\textwidth}@{}}
	\toprule
	\textbf{Status} & \textbf{Meaning} \\
	\midrule
	\texttt{Running} & Executing, or waiting to be resumed. \\
	\texttt{Suspended} & Stopped at a resume key --- an approval, a timer --- and resumable. \\
	\texttt{Succeeded} & Reached its terminal node and reported success. \\
	\texttt{Failed} & Reached a terminal node and reported failure, or was cancelled. \\
	\texttt{Refused} & An agent call was refused before any spawn. It would not start,
	which is not the same as failing. \\
	\bottomrule
\end{tabular}
\end{center}

\begin{center}
\resizebox{\linewidth}{!}{
\begin{tikzpicture}[node distance=14mm and 18mm]
  \node[b] (start) {run start\\\scriptsize an id, then a journal};
  \node[v, right=of start] (live) {live?};
  \node[b, right=of live, yshift=14mm] (susp) {Suspended\\\scriptsize resumable with a key};
  \node[b, right=of live, yshift=-14mm] (term) {terminal\\\scriptsize Succeeded or Failed};
  \node[b, below=of start] (refused) {Refused\\\scriptsize refused before spawn};
  \draw[f] (start) -- (live);
  \draw[f] (live) -- node[p,above=2pt,midway]{needs a key} (susp);
  \draw[f] (live) -- node[p,below=2pt,midway]{ran} (term);
  \draw[sf] (susp.north) to[bend left=40] node[p,above=2pt,midway]{resumed} (live.north);
  \draw[bk] (live) -- node[p,left=2pt,midway]{cannot start} (refused);
\end{tikzpicture}}
\captionof{figure}{Five statuses, and the loop that matters: a suspended run is not a finished
one. Retention keeps a Running or Suspended run however old it is, because the daemon may still
resume it.}
\end{center}

A direct run reports one of three terminals, and \texttt{pii graph run} maps them to exit
codes an operator can branch on: success is 0, a failure is 1, a suspended run is 4, and a
graph the validator refuses is 3 --- before anything is executed and, deliberately, before
any store file is created.

\section{What it remembers, and where}

Seven store files appear in \texttt{pii health}. Six of them are what a backup copies:
the main store, memory, approvals, heal, channel and timer. The queue has its own file
and is reported alongside them, but it is not in the backup set at this revision --- a
detail that matters the first time you restore and find the work list empty.

\begin{center}
\resizebox{\linewidth}{!}{
\begin{tikzpicture}[node distance=8mm and 12mm]
  \node[s] (main) {store.redb\\\scriptsize runs, journal, sessions};
  \node[s, right=of main] (mem) {memory.redb\\\scriptsize documents, entities};
  \node[s, right=of mem] (aux) {approval, heal\\\scriptsize decisions and gates};
  \node[s, below=of main] (q) {queue.redb\\\scriptsize work rows};
  \node[s, right=of q] (dur) {durable/\\\scriptsize channel \textperiodcentered{} timer};
  \node[b, below=of q, xshift=22mm] (dir) {state dir\\\scriptsize \texttt{PII\_STATE\_DIR} or
\texttt{\$HOME/.local/state/pii}};
  \draw[sf] (main) -- (dir);
  \draw[sf] (mem) -- (dir);
  \draw[sf] (aux) -- (dir);
  \draw[sf] (q) -- (dir);
  \draw[sf] (dur) -- (dir);
\end{tikzpicture}}
\captionof{figure}{Every store except the socket lives under one state directory. The durable
pair sits in \texttt{durable/} because the runtime directory is a tmpfs on a normal host, and a
channel cursor or a timer that vanishes at reboot is worse than no timer at all.}
\end{center}

The main store is a redb file, which never shrinks on its own: space is returned only at
the end of the file, and only by compaction. That is the reason a retention policy exists
at all, and the reason \texttt{pii store compact} is documented as narrow --- it reclaims
free pages inside the file, not live history.

\section{Capabilities are a ceiling}

Twelve capability bits exist, and every one of them is checked twice: the graph declares
what it may hold, and each node may hold only what its graph holds. A node that asks for
more is refused at load, not degraded at runtime.

\begin{center}
\resizebox{\linewidth}{!}{
\begin{tikzpicture}[node distance=10mm and 12mm]
  \node[b] (ceiling) {operator ceiling\\\scriptsize \texttt{grant\_ceiling}};
  \node[b, right=of ceiling] (graph) {graph \texttt{caps}\\\scriptsize in the file};
  \node[b, right=of graph] (node) {node \texttt{caps}\\\scriptsize the intersection};
  \node[v, right=of node] (allow) {permitted?};
  \node[b, below left=of allow] (yes) {the node runs};
  \node[b, below right=of allow] (no) {refused at load\\\scriptsize named capability};
  \draw[f] (ceiling) -- (graph);
  \draw[f] (graph) -- (node);
  \draw[f] (node) -- (allow);
  \draw[f] (allow) -- node[p,left=2pt,midway]{yes} (yes);
  \draw[bk] (allow) -- node[p,right=2pt,midway]{no} (no);
\end{tikzpicture}}
\captionof{figure}{The ladder narrows downwards and never widens. A \texttt{call} can grant only
what its host grants, which is why \texttt{night} holds what \texttt{queue} needs before it may
call it.}
\end{center}

\begin{longtable}{@{}p{0.20\textwidth}p{0.72\textwidth}@{}}
	\toprule
	\textbf{Capability} & \textbf{What holding it allows} \\
	\midrule
	\endhead
	\texttt{fs\_read} & Read files inside an allow path scope. \\
	\texttt{fs\_write} & Write files inside an allow path scope. \\
	\texttt{proc} & Run commands --- this is what lets a judge execute exit criteria. \\
	\texttt{spawn} & Start an agent or a child process. Every model call needs it. \\
	\texttt{net} & Reach the network, subject to the config's host grants. \\
	\texttt{hw} & Touch hardware through the HAL. \\
	\texttt{secrets} & Read secret material. \\
	\texttt{memory\_read}, \texttt{memory\_write} & Read from, or write into, the agent's
	memory. \\
	\texttt{post} & Send a message into a configured channel. \\
	\texttt{mqtt\_read}, \texttt{mqtt\_write} & List and subscribe to a broker's devices;
	command one through its allowlist. \\
	\bottomrule
\end{longtable}

The measured consequence is worth stating plainly: a coder inside a graph whose
\texttt{caps} omitted \texttt{fs\_write} did not silently write nothing. It stopped and
reported \enquote{no write or execute tool is available}. That sentence is the feature
working; a degraded coder that pretends to work is the failure it replaced.

\section{How a run is named, and what a trace looks like}

Run identifiers are allocated by the store, not by the caller, so two processes cannot
collide: the daemon's own logs and the graph files quote them in the form
\texttt{r000000000130}. A caller may pass \texttt{--run-id} on the offline path, which is
how the test suites keep a fixture's journal findable.

\texttt{pii health} is the widest single view of what a process is. Its text form is
short enough to quote in a bug report, and every line is a fact a command can be held to:

\begin{lstlisting}[language=bash]
pii <version> (<revision>, <target>)
client: <revision>
config: <path> sha256=<hash>
stores:
  main: <path> (<bytes> bytes, format_version=<n>)
  memory: <path> (<bytes> bytes, format_version=<n>)
mounts:
  workspace [rw]: <root>  (<policy>)
since_deploy:
  starts: <n>
counts:
  runs: <n>
\end{lstlisting}

The field names, their order and their punctuation are exact; the values are whatever your
daemon reports. Read it as four claims, in order. The first line is \emph{what the daemon
is}: version, revision and target. The second is \emph{what your client is}: if the two
revisions differ, \texttt{pii health} exits 3 and prints a line beginning \texttt{DRIFT},
and that is the correct reading immediately after installing a new binary and before
restarting the service. The \texttt{config} line is the hash of the file the daemon actually
loaded --- not the path you think it loaded. The \texttt{stores} block is the one that
catches the most common installation mistake: a daemon whose store list is empty is a
daemon pointed at an empty directory, and it will happily report itself active.

\section{What pii refuses to claim}

Three refusals are part of the design, and an administrator meets all three.

\begin{itemize}
	\item A store whose format version is newer than the binary supports is not read
	half-way. The daemon exits 2 and writes nothing to the file, with the version and the
	path in the message.
	\item A health answer whose build revision is unknown on either side is reported as
	unverified, with a non-zero exit code, rather than as agreement. Unknown is never
	agreement.
	\item A retention sweep will not remove a row it cannot date. A run created before the
	clock was recorded, or a fixture with no timestamp, is kept and counted, because
	deleting on a missing timestamp destroys exactly the runs an operator is looking for.
\end{itemize}

\section{What it deliberately is not}

Four absences are decisions rather than gaps, and each one shapes how you operate the
thing.

\begin{itemize}
	\item \textbf{Not a cluster.} The stores are redb files with an exclusive lock, and the
	daemon claims the durable pair before it binds its socket. A second daemon pointed at
	the same state directory is refused, by name, at start. Scaling out means changing the
	store, not adding a replica.
	\item \textbf{Not an online compactor.} \texttt{pii store compact} runs redb's in-place
	compaction, and it needs the daemon stopped, because it takes the same exclusive lock
	the daemon holds for its lifetime. There is no background compaction to tune or to
	fear.
	\item \textbf{Not a cron daemon.} A \texttt{[[cron]]} entry carries a
	\texttt{schedule} string, but that string is a display mirror: what fires a graph is a
	\texttt{timer} node at its entry. A graph whose entry is not a timer will not be
	started by the schedule, and the config will look correct while nothing runs.
	\item \textbf{Not a TLS client for MQTT.} The broker client speaks plain TCP. A remote
	broker with credentials needs an explicit acknowledgement in the config, and the
	honest answer for a non-loopback broker is a tunnel.
\end{itemize}

Knowing these four saves the two most expensive mistakes in operation: pointing a second
process at a live state directory, and believing a schedule that only mirrors itself.

\section{How to read this manual}

Chapter~\ref{ch:install} gets a host running: the two directories, the installer, the
unit, the credentials, and the three commands that prove the result. Chapter~\ref{ch:daily}
is what a normal day looks like --- health, one-shot prompts, graphs by name and by file,
the queue, and how to read a trace. Chapter~\ref{ch:operate} is the part nobody reads
until something breaks: backup and restore, retention, upgrades, logs, disk ownership, the
deploy card and the failure modes that have actually happened.

Chapters~\ref{ch:commands} and~\ref{ch:config} are reference material and are meant to be
searched, not read. Chapter~\ref{ch:observability} is the metrics and event surface, and
Chapter~\ref{ch:graphs} is the graph library: which graphs an operator runs and what each
will do to the machine.

Two conventions run through the book. A path in \texttt{typewriter} is literal and was
copied from the repository or from a command's own output. A number in the text is a
measurement recorded in the repository, with the row or document that recorded it named
where the number is load-bearing. Where this manual and a design document disagree, the
manual is describing the code at the built revision, and the disagreement is itself worth
reporting.
